%! Exponential Function Context and Data Modeling — Unit 4, Lesson 4.6 — Guided Notes (Answer Key)
\documentclass[10pt]{article}
\usepackage{precalculus-article}
\usepackage{precalculus-key}   % pulls in -boxes; do NOT also load -boxes
\newcommand{\TallMath}[1]{$\displaystyle #1\rule[-1.4em]{0pt}{3.2em}$}

\begin{document}

\pageheader{Unit 4, Lesson 4.6}{Guided Notes --- Answer Key}
\namedateperiod

\vspace{0.08in}

\begin{objectivebox}
By the end of this lesson, I will be able to\ldots
\begin{enumerate}[itemsep=2pt]
  \item Construct an exponential model from a \ans{growth} factor and \ans{initial} value,
        or from \ans{two} input-output pairs. \hfill\textit{(LO 2.5.A)}
  \item Interpret the base $b$ as a \ans{growth} factor and connect it to \ans{percent} change. \hfill\textit{(LO 2.5.B)}
  \item Use the model to \ans{predict} values within a reasonable \ans{domain}. \hfill\textit{(LO 2.5.B)}
\end{enumerate}
\end{objectivebox}

\vspace{0.06in}

\begin{vocabbox}
\termblanklong{exponential model}
\termblanklong{growth factor}
\termblanklong{percent change}
\termblanklong{initial value}
\termblanklong{natural base $e$}
\termblanklong{prediction}
\termblanklong{domain restriction}
\end{vocabbox}

\vspace{0.06in}

\begin{hookbox}
A town's population over four years:
\begin{center}
\renewcommand{\arraystretch}{1.4}
\begin{tabular}{c|c|c}
\textbf{Year} & \textbf{Population} & \textbf{Ratio (successive)} \\ \hline
0 & 10{,}000 & --- \\
1 & 10{,}300 & $10300 \div 10000 = $ \ans{1.03} \\
2 & 10{,}609 & $10609 \div 10300 \approx $ \ans{1.03} \\
3 & 10{,}927 & $10927 \div 10609 \approx $ \ans{1.03} \\
\end{tabular}
\end{center}
\medskip
The ratios are \ans{constant (equal)}, so the growth is \ans{exponential}. The growth rate per year is \ans{3\%}.
\end{hookbox}

\vspace{0.06in}

\begin{notesbox}{LO 2.5.A --- Building a Model from Initial Value and Ratio}
\textbf{General Form:} $f(x) = a \cdot b^x$, where $a = $ \ans{initial value} and $b = $ \ans{growth factor}.

\medskip
\textbf{Connecting $b$ to percent change (EK 2.5.B.1):}
\begin{itemize}[itemsep=3pt]
  \item Growth at rate $r$: $b = $ \ans{$1 + r$}
  \item Decay at rate $r$: $b = $ \ans{$1 - r$}
\end{itemize}

\medskip
\textbf{Example 1:} A population starts at 10{,}000 and grows 3\% per year. Write the model.

\smallskip
$a = $ \ans{10{,}000} \quad $b = $ \ans{1.03} \quad Model: $P(t) = $ \ans{$10000 \cdot (1.03)^t$}

\medskip
\textbf{Example 2:} A car's value starts at \$28{,}000 and depreciates 12\% per year. Write the model.

\smallskip
$a = $ \ans{28{,}000} \quad $b = $ \ans{0.88} \quad Model: $V(t) = $ \ans{$28000 \cdot (0.88)^t$}

\medskip
\textbf{Example 3:} A bank account earns 4.5\% interest annually on an initial \$1{,}000.

\smallskip
Model: $A(t) = $ \ans{$1000 \cdot (1.045)^t$} \quad After 6 years: $A(6) = $ \ans{$\approx \$1308.65$}
\end{notesbox}

\vspace{0.06in}

\begin{notesbox}{LO 2.5.A --- Building a Model from Two Points (EK 2.5.A.3)}
\textbf{Strategy:} Given $(x_1, y_1)$ and $(x_2, y_2)$, write two equations:
\[
  y_1 = a \cdot b^{x_1} \qquad \text{and} \qquad y_2 = a \cdot b^{x_2}
\]
Divide the second by the first to eliminate $a$:
\[
  \frac{y_2}{y_1} = b^{\,x_2 - x_1} \qquad \Rightarrow \qquad b = \left(\frac{y_2}{y_1}\right)^{1/(x_2-x_1)}
\]
Then solve for $a$: $a = y_1 / b^{x_1}$.

\medskip
\textbf{Example:} An exponential function passes through $(0, 5)$ and $(3, 40)$.

\smallskip
Step 1 --- Find $b$: \ans{$b = (40/5)^{1/3} = 8^{1/3} = 2$}

\smallskip
Step 2 --- Find $a$: \ans{$a = 5 / 2^0 = 5$}

\smallskip
Model: $f(x) = $ \ans{$5 \cdot 2^x$}

\medskip
\textbf{Natural base $e$ (EK 2.5.A.6):} $e \approx $ \ans{2.718}. Used in continuous models: $f(t) = a \cdot e^{rt}$.

Example: $f(t) = 500 e^{0.04t}$ grows \ans{4}\% continuously per year.
\end{notesbox}

\vspace{0.06in}

\begin{notesbox}{LO 2.5.B --- Equivalent Forms and Predictions (EK 2.5.B.2)}
\textbf{Equivalent forms reveal different time-scale properties.}

$f(d) = 3 \cdot 2^d$ models growth where $d$ = days.

Rewrite for weeks (let $w = d/7$, so $d = 7w$):
\[
  f(d) = 3 \cdot 2^d = 3 \cdot 2^{7w} = 3 \cdot (2^7)^w = 3 \cdot \ans{128}^w
\]
The base \ans{128} means the quantity multiplies by \ans{128 (i.e., increases by a factor of 128)} each week.

\medskip
\textbf{Prediction and domain caution:} Use the model for inputs \ans{within the realistic context}.

Predict the town population at $t = 20$: $P(20) = 10000 \cdot (1.03)^{20} \approx $ \ans{18{,}061}

\textit{Domain note:} Is $t = 200$ realistic? \ans{Probably not; exponential growth models break down far from the data.}
\end{notesbox}

\begin{teachernote}
Hook table: each ratio is 1.03, confirming 3\% per-year exponential growth.

Example 3: $1000(1.045)^6 = 1000 \times 1.30865 \approx \$1308.65$.

Two-point example: emphasize the division step. Students often try to subtract equations instead.

Equivalent forms: $2^7 = 128$. The new base tells you the weekly multiplier, not a new rate.

Domain caution: model validity depends on the context. Exponential growth cannot continue indefinitely.
\end{teachernote}

\end{document}
